\section{Introduction}
\label{sec:introduction}

Wi-Fi sensing infers human motion from variations of channel state information
(CSI), reusing commodity radios without recording identifiable images or audio.
Its practical value, however, depends on representations that separate motion
from environment-, device-, and packet-dependent distortions. SHARP addresses
this requirement by transforming CSI into phase-stable Doppler maps before
classification, and reports person- and environment-independent human activity
recognition (HAR) with commodity access points \cite{meneghello2023sharp}.
The transformation is not a plain Fourier transform: it estimates sparse paths,
selects a phase reference, sanitizes nuisance phase, and only then applies the
Doppler analysis.

This project asks whether the expensive part of that pipeline can be
approximated directly. Given a local window of raw complex CSI, a neural network
is trained to reproduce the Doppler map generated by SHARP. If successful, this
would preserve an established classifier interface while replacing iterative
signal processing with a fixed feed-forward computation. It also poses a useful
scientific test: can generic convolutional networks learn the phase invariance
that model-based preprocessing imposes explicitly?

The primary research question is:
\emph{Can an end-to-end neural network replace SHARP's optimization-intensive
CSI preprocessing while preserving both sparse motion structure and downstream
HAR performance?} Three supporting questions separate possible causes of
failure. First, do frequency-aware decoding, antenna-independent weights, exact
temporal alignment, and motion-aware training prevent a center-spectrum
shortcut? Second, do pixel reconstruction metrics predict whether generated
maps remain useful to a frozen SHARP classifier? Third, what does a comparison
with explicit affine phase correction reveal about the transformation that a
student must learn?

The contribution is a feasibility study rather than a new state-of-the-art HAR
method:
\begin{itemize}
    \item a reproducible CSI-to-SHARP-map distillation pipeline with controlled
    chronological splits and task-aware evaluation;
    \item a progressively corrected architecture and training setup addressing
    antenna sharing, subcarrier structure, temporal support, sparse motion, and
    Doppler locality;
    \item evidence that a full-resolution student partially preserves held-out
    motion and HAR utility, together with a class-level analysis of the
    remaining smoothing failure and a stronger explicit affine baseline.
\end{itemize}

\figref{fig:pipeline} summarizes the comparison. Both pixel fidelity and
classifier utility use the official precomputed SHARP representation: students
regress its maps, and each generated representation is sent through a classifier
trained on the same map distribution. The rest of the report reviews related
representations, defines the data and learning problem, presents the final model
and loss, and answers the research questions experimentally.

\begin{figure*}[t]
    \centering
    \includegraphics[width=\textwidth]{pipeline.pdf}
    \caption{Evaluation design. Raw CSI is transformed by the SHARP teacher, a
    direct neural student, or a simple affine phase correction followed by the
    same fixed Hann/STFT tail. Reconstruction metrics compare maps with the
    official SHARP target; a frozen classifier measures retained HAR utility.}
    \label{fig:pipeline}
\end{figure*}
